\clearpage\pagestyle{otMain}\setlength{\otBlocEspacement}{2.5mm}\fontsize{11}{15}\selectfont\corrigesChapitre{2}{Multiples, diviseurs et partages}
\begin{blocLivre}{cor-2-1}\begin{corrige}{1}{Quotient et reste}\demarche{Encadrer par deux multiples.}
$8\times8=64$ et $67-64=3$. Donc $\res{67=8\times8+3}$, avec $0\le3<8$.
\end{corrige}\end{blocLivre}
\begin{blocLivre}{cor-2-2}\begin{corrige}{2}{Diviseur ou multiple}\demarche{Écrire un produit entier.}
$54=6\times9$. 54 est un $\res{\text{multiple}}$ de 6 et 6 un $\res{\text{diviseur}}$ de 54.
\end{corrige}\end{blocLivre}
\begin{blocLivre}{cor-2-3}\begin{corrige}{3}{Tester 450}\demarche{Choisir le bon critère.}
Il finit par 0 : divisibilité par 2, 5 et 10. La somme des chiffres vaut 9 : divisibilité par 3 et 9. Les cinq réponses sont $\res{\text{oui}}$.
\end{corrige}\end{blocLivre}
\begin{blocLivre}{cor-2-4}\begin{corrige}{4}{Lister des diviseurs}\demarche{Former des paires de facteurs.}
Les paires sont $1\times18$, $2\times9$, $3\times6$. La liste est $\res{1,2,3,6,9,18}$.
\end{corrige}\end{blocLivre}
\begin{blocLivre}{cor-2-5}\begin{corrige}{5}{Valider une écriture}\demarche{Contrôler aussi la taille du reste.}
Non : le reste 11 est supérieur au diviseur 6. Il faut $\res{59=6\times9+5}$ ; cette fois $0\le5<6$.
\end{corrige}\end{blocLivre}
\begin{blocLivre}{cor-2-6}\begin{corrige}{6}{Un chiffre manquant}\demarche{Tester les dix chiffres avec le critère.}
La somme des chiffres est $5+a$. Elle peut valoir 6, 9 ou 12 : $\res{a=1,4,7}$. On retrouve 231, 234 et 237.
\end{corrige}\end{blocLivre}
\begin{blocLivre}{cor-2-7}\begin{corrige}{7}{Deux critères ensemble}\demarche{Croiser les critères.}
$7+a+5=12+a$ doit être multiple de 9 : seul 18 convient, donc $\res{a=6}$. Le nombre 765 finit par 5 : il est aussi divisible par 5.
\end{corrige}\end{blocLivre}
\begin{blocLivre}{cor-2-8}\begin{corrige}{8}{Diviseurs communs}\demarche{Filtrer une première liste.}
Diviseurs de 24 : $1,2,3,4,6,8,12,24$. Ceux qui divisent aussi 36 sont $\res{1,2,3,4,6,12}$.
\end{corrige}\end{blocLivre}
\begin{blocLivre}{cor-2-9}\begin{corrige}{9}{Premier multiple commun}\demarche{Comparer les listes de multiples.}
Multiples de 6 : $6,12,18,24,\ldots$ ; de 8 : $8,16,24,\ldots$. Le premier commun est $\res{24}$.
\end{corrige}\end{blocLivre}
\begin{blocLivre}{cor-2-10}\begin{corrige}{10}{Un crible limité}\demarche{Barrer les multiples organisés.}
On retire 1, puis les multiples de 2 et de 3 autres que 2 et 3. Les nombres restants sont $\res{2,3,5,7,11,13,17,19}$. Les nombres composés de cet intervalle ont un facteur 2 ou 3.
\end{corrige}\end{blocLivre}
\begin{blocLivre}{cor-2-11}\begin{corrige}{11}{Transporter une classe}\demarche{Arrondir dans le sens du besoin.}
$137=9\times15+2$. Quinze minibus ne suffisent pas : il en faut $\res{16}$. Ils offrent 144 places, donc $\res7$ places libres.
\end{corrige}\end{blocLivre}
\begin{blocLivre}{cor-2-12}\begin{corrige}{12}{Des bouquets identiques}\demarche{Écrire la contrainte de divisibilité.}
Les diviseurs communs sont $1,2,7,14$. Au moins trois bouquets impose $\res{7\text{ ou }14}$. Avec 14 bouquets : $\res3$ tulipes et $\res4$ roses chacun.
\end{corrige}\end{blocLivre}
\begin{blocLivre}{cor-2-13}\begin{corrige}{13}{Deux signaux}\demarche{Chercher un multiple commun minimal.}
Les premiers multiples communs apparaissent à 36 minutes : $3\times12=2\times18=36$. Le prochain déclenchement commun est à $\res{9\ \mathrm h\ 36}$.
\end{corrige}\end{blocLivre}
\begin{blocLivre}{cor-2-14}\begin{corrige}{14}{Un nombre à retrouver}\demarche{Appliquer les contraintes successivement.}
Les entiers de la forme $7q+2$ dans cet intervalle sont 44, 51 et 58. Seul $\res{51}$ est divisible par 3. Contrôle : $51=7\times7+2$.
\end{corrige}\end{blocLivre}
\begin{blocLivre}{cor-2-15}\begin{corrige}{15}{Une affirmation générale}\demarche{Expliquer ce qui reste vrai pour tous les cas.}
Chaque nombre est formé de paquets de quatre ; réunir les deux collections donne encore des paquets de quatre. Ici $28+44=\res{72=4\times18}$. L’argument des paquets vaut pour tous les multiples, au-delà de cet exemple.
\end{corrige}\end{blocLivre}
